\subsubsection{Gradient Episodic Memory (GEM)}

\gls{gem} \cite{ll_gem_2017} is an experience replay method that mitigates \gls{cf}
by constraining the direction of gradient updates during training on a new task.
Rather than penalising weight changes directly, as in regularisation-based
approaches such as \cite{ll_ewc_2017}, \gls{gem} enforces that the loss on
previously observed tasks does not increase as a result of any update to the
model parameters. \gls{gem} requires task identity during training to index the
per-task episodic memory buffers, but does not require task identity at
inference time. It is therefore compatible with both \gls{til} and \gls{cil}
evaluation settings, provided a shared output head is used in the \gls{cil}
case.


\paragraph{Episodic Memory and Gradient Constraints}

\gls{gem} maintains a fixed-size episodic memory buffer $\mathcal{M}_k$ for
each observed task $k$, populated with a subset of samples encountered during
that task's training phase. At each subsequent training step on task $t$, the
gradient of the current loss, $g = \nabla_\theta \mathcal{L}^t$, is computed
alongside reference gradients $g_k = \nabla_\theta \mathcal{L}^{\mathcal{M}_k}$
for all prior tasks $k < t$. A gradient update is considered safe with respect
to task $k$ if it does not increase the episodic loss for that task. Using the
dot product to test orthogonality (Section~\ref{sect_bg_la}), this is expressed
as:
%
\begin{equation}\label{eqn:gem_constraint}
    \langle g, g_k \rangle \geq 0, \quad \forall k < t
\end{equation}
%
If the proposed gradient $g$ violates any of these constraints, it is projected
onto the feasible region defined by \eqref{eqn:gem_constraint}. Following the
orthogonal projection formulation in Section~\ref{sect_bg_la}, this is achieved
by solving a \gls{qp} that finds the minimal correction to $g$:
%
\begin{equation}\label{eqn:gem_qp}
    \tilde{g} = \argmin_{z} \frac{1}{2}\|z - g\|^2_2
    \quad \text{subject to} \quad Gz \geq 0
\end{equation}
%
where $G = [g_1, g_2, \ldots, g_{t-1}]^\top$ collects the reference gradients
from all prior tasks. Geometrically, $\tilde{g}$ is the projection of $g$ onto
the intersection of half-spaces defined by the per-task constraints; it
eliminates the components of $g$ that point away from prior task optima whilst
preserving the remaining update direction. If no constraint is violated, $g$ is
used directly without modification.

\paragraph{Complexity Considerations}

The \gls{qp}  in \eqref{eqn:gem_qp} grows in both size and computational cost as
the number of tasks increases, since one constraint is added per task. This
makes \gls{gem} computationally expensive in long task sequences, motivating the \gls{agem} approach (Section~\ref{sec:agem}).

